\documentclass[11pt]{article}
\usepackage[a4paper,margin=26mm]{geometry}
\usepackage[T1]{fontenc}
\usepackage{lmodern}
\usepackage{amsmath,amssymb,amsthm}
\usepackage{microtype}
\usepackage{array}
\usepackage{xurl}
\usepackage{hyperref}
\hypersetup{colorlinks=true,linkcolor=blue,urlcolor=blue,citecolor=blue,
 pdftitle={A non-hollow lattice 4-simplex of covering radius 2 outside the equality list of a conjecture of Codenotti, Santos and Schymura},
 pdfauthor={Horace Dong}}
\newtheorem{theorem}{Theorem}
\newtheorem{lemma}{Lemma}
\newtheorem{proposition}{Proposition}
\newtheorem{corollary}{Corollary}
\theoremstyle{definition}
\newtheorem{definition}{Definition}
\theoremstyle{remark}
\newtheorem{remark}{Remark}
\newcommand{\repoversion}{1.25.0}
\setlength{\emergencystretch}{1.5em}
\newcommand{\R}{\mathbb R}
\newcommand{\Z}{\mathbb Z}
\newcommand{\one}{\mathbf 1}
\newcommand{\conv}{\operatorname{conv}}
\newcommand{\Vol}{\operatorname{Vol}}
\newcommand{\vol}{\operatorname{vol}}
\newcommand{\G}{\mathcal G}
\newcommand{\tot}{\operatorname{tot}}

\title{A non-hollow lattice $4$-simplex of covering radius $2$\\ outside the
equality list of a conjecture of\\ Codenotti, Santos and Schymura}
\author{Horace Dong\\ \small Independent researcher\\
\small\texttt{maxwelldhx@gmail.com}\\
\small ORCID: \href{https://orcid.org/0009-0002-7554-7548}{0009-0002-7554-7548}}
\date{22 November 2026}
\begin{document}
\maketitle

\begin{abstract}
The covering radius $\mu(P)$ of a convex body $P\subseteq\R^d$ is the least
$\mu\ge0$ such that the lattice translates of $\mu P$ cover $\R^d$. Codenotti,
Santos and Schymura conjectured that every non-hollow lattice $d$-polytope $P$
satisfies $\mu(P)\le d/2$, with equality if and only if $P$ is obtained by
direct sums and/or translations of the standard terminal simplices
$S(\one_l)$ (their Conjecture~A); they proved this for $d\le3$. We show that
the $4$-simplex $P=\conv\{e_1,e_2,e_3,e_4,-(2,1,1,1)\}$, whose only interior
lattice point is the origin, has covering radius exactly $2=d/2$ and
normalized volume~$6$, and that it is not obtained from the simplices
$S(\one_l)$ by direct sums and lattice translations. So the inequality of
Conjecture~A holds with equality for $P$, while the ``only if'' part of its
equality clause fails. The simplex itself is known from toric geometry.
Lean~4 verifies that the origin is an interior point of $P$, that
$\mu(P)=2$, that the normalized volume is~$6$, and that $P$ lies outside a
family which, by a short argument on the definitions given in the text,
contains all direct sums and lattice translations of the simplices
$S(\one_l)$. The proofs were found and formalized by AI agents working under
the author's direction, as described at the end of the paper.
\end{abstract}

\section{Introduction}

Let $K\subseteq\R^d$ be a convex body and let $\Lambda$ be a lattice. Codenotti,
Santos and Schymura~\cite{css} write the covering radius as
\[
\mu(K,\Lambda)=\min\{\mu\ge0:\ \mu K+\Lambda=\R^d\},
\]
and $\mu(K)=\mu(K,\Z^d)$. A lattice polytope is a polytope whose vertices are
lattice points; it is \emph{non-hollow} if its interior contains a lattice
point. The covering radius is invariant under unimodular transformations
$x\mapsto Ux+z$ ($U\in GL_d(\Z)$, $z\in\Z^d$)~\cite[p.~2]{css}. For
$\omega\in\R^{d+1}_{>0}$ the authors define
$S(\omega)=\conv\{-\omega_0\one_d,\omega_1e_1,\dots,\omega_de_d\}$, where
$\one_d=(1,\dots,1)$. The \emph{standard terminal simplex} is
$S(\one_{d+1})=\conv\{-\one_d,e_1,\dots,e_d\}$, and
$\mu(S(\one_{d+1}))=d/2$~\cite[Eq.~(1)]{css}. If $K\subseteq V$ and
$L\subseteq W$ are convex bodies containing the origin in complementary linear
subspaces $V$ and $W$ of $\R^d$, their direct sum is
$K\oplus L=\{\lambda x+(1-\lambda)y: x\in K,\ y\in L,\ \lambda\in[0,1]\}$,
taken with the direct sum of the two lattices~\cite[Section~2.1]{css}; the
covering radius is additive on direct sums~\cite[Corollary~2.2]{css}. Hence,
as the authors note, ``direct sums of simplices of the form $S(\one_l)$ or
lattice translates thereof'' have covering radius $d/2$~\cite[p.~2]{css}. They
conjectured:

\begin{quote}
\textbf{Conjecture A of~\cite{css}.} Let $P\subseteq\R^d$ be a non-hollow
lattice $d$-polytope. Then $\mu(P)\le d/2$, with equality if and only if $P$
is obtained by direct sums and/or translations of simplices of the form
$S(\one_l)$.
\end{quote}

\noindent The wording is the same in arXiv:1903.02866v2 and in the journal
version; page numbers refer to arXiv:1903.02866v2. The authors prove
Conjecture~A for $d\le3$~\cite[Theorem~1.3]{css}; in dimension three the
equality cases are nine polytopes up to unimodular
equivalence~\cite[Example~1.1 and Lemma~3.8]{css}. Example~1.1 of~\cite{css}
translates a factor before the direct sum, so that the origin can lie on the
boundary of a factor. A later restatement reads the equality list up to
unimodular transformation, as ``direct sums of translates of standard terminal
simplices''~\cite[p.~2 and Conjecture~4.1]{kr}.

Throughout, $u=(2,1,1,1)$ and
\[
P=\conv\{e_1,e_2,e_3,e_4,-u\}\subseteq\R^4,
\]
with the lattice $\Z^4$. We write $\Vol(K)$ for the normalized volume, that is
the Euclidean volume divided by the volume of a unimodular simplex, as
in~\cite[p.~3]{css}.

\begin{theorem}\label{thm:main}
\begin{enumerate}
\item[\textup{(a)}] The origin lies in the interior of $P$; so $P$ is a
non-hollow lattice $4$-polytope.
\item[\textup{(b)}] $\mu(P)=2$: the lattice translates of $2P$ cover $\R^4$,
and for every $0\le r<2$ some point of $\R^4$ lies in no lattice translate of
$rP$.
\item[\textup{(c)}] $\Vol(P)=6$, that is, $\vol(P)=6\,\vol(\conv\{0,e_1,e_2,e_3,e_4\})$.
\item[\textup{(d)}] $P$ does not belong to the family $\G$ of
Definition~\ref{def:G}.
\end{enumerate}
\end{theorem}

\noindent All four parts are formally verified in Lean
(Section~\ref{sec:lean}). The family $\G$ is generated from finite sets of at
least two lattice points with sum zero by lattice translations and by convex
hulls of unions. Lemma~\ref{lem:G} shows that $\G$ contains every polytope
obtained from simplices $S(\one_l)$ by direct sums and lattice translations,
and every unimodular image of such a polytope. Lemma~\ref{lem:G} is an
argument about the definitions in~\cite{css}; it is not a Lean theorem,
because the Lean files do not define direct sums. With it,
Theorem~\ref{thm:main} gives:

\begin{corollary}\label{cor:A}
$P$ is a non-hollow lattice $4$-polytope with $\mu(P)=2=d/2$ that is not
obtained, even up to a unimodular transformation, by direct sums and/or
lattice translations of simplices $S(\one_l)$. Every unimodular image of $P$
has the same properties. So the ``only if'' direction of the equality clause
in Conjecture~A of~\cite{css} fails in dimension~$4$, the smallest dimension
not covered by~\cite[Theorem~1.3]{css}, when ``translations'' are read as
lattice translations, as on p.~2 of~\cite{css}.
\end{corollary}

\noindent The inequality $\mu(P)\le d/2$ of Conjecture~A holds for $P$, with
equality; the example concerns only the equality clause. If arbitrary real
translations are allowed in the equality clause, $P$ is still not in the
list, by a volume argument that is not formalized (Remark~\ref{rem:real}).
The same example also contradicts the equality clause for $k=1$ of
Conjecture~D of~\cite{css}, and it does not refute their Conjecture~C
(Section~\ref{sec:remarks}); neither statement is in Lean.

\subsection*{Prior work}
The simplex $P$ is not new. Its vertices are primitive, they span $\Z^4$, and
they satisfy $2e_1+e_2+e_3+e_4+(-u)=0$; so $P$ is, up to a unimodular
transformation, the simplex spanned by the primitive ray generators of the
fan of the weighted projective space $\mathbb P(1,1,1,1,2)$. Kasprzyk's thesis
shows that a Gorenstein Fano weighted projective space of dimension four with
at worst terminal singularities is $\mathbb P^4$ or
$\mathbb P(1,1,1,1,2)$~\cite[Proposition~9.7.10]{kas}. Lee~\cite[Appendix~C]{lee} uses
the simplex $\conv\{e_1,e_2,e_3,e_4,(-1,-1,-1,-2)\}$, which is $P$ after a
permutation of the coordinates, as an example in mirror symmetry; he lists
its lattice points and facet normals and notes that its normalized volume is
six. The value $\Vol(P)=6$ is therefore known; we include it because it
separates $P$ from the equality list when arbitrary translations are allowed.

Codenotti, Santos and Schymura give a closed formula for the covering radius
of the simplices $S(\omega)$~\cite[Theorem~1.4]{css}. It does not apply to $P$,
because $P$ is not a unimodular image of any $S(\omega)$
(Remark~\ref{rem:somega}). Their Appendix~B computes covering radii of lattice
simplices from the diameter of a Cayley digraph~\cite[Lemmas~B.1
and~B.2]{css}; for $P$ this is a digraph on $216$ vertices, and our proof of
$\mu(P)\le2$ is a short certificate of the same kind (Remark~\ref{rem:cayley}).
As far as we know, the covering radius of $P$ and its relation to the
equality clause of Conjecture~A have not been stated before; the searches are
described in the final section. Recent work on related questions includes
Krivokuća's bounds on covering minima~\cite{kr} and the paper of Codenotti,
Freyer and Krivokuća on minimal covering bodies~\cite{cfk}.

\section{Coordinates and lattice points}\label{sec:coord}

For $x\in\R^4$ write $\tot(x)=x_1+x_2+x_3+x_4$, and let $\one=(1,1,1,1)$. Let
\[
M=I+u\one^{\mathsf T}=\begin{pmatrix}3&2&2&2\\1&2&1&1\\1&1&2&1\\1&1&1&2\end{pmatrix},
\qquad \det M=1+\tot(u)=6 .
\]
Since $\tot(u)=5$ we have $Mu=6u$ and $\tot(Mb)=6\tot(b)$. For $r\in\R$ define the
affine map $X_r(b)=Mb-ru$. Let $\Delta=\{b\in\R^4: b_i\ge0,\ \tot(b)\le1\}$ be the
standard simplex $\conv\{0,e_1,\dots,e_4\}$.

\begin{lemma}\label{lem:coord}
\begin{enumerate}
\item[\textup{(i)}] $X_r$ is a bijection of $\R^4$ (since $\det M=6\ne0$). Its inverse sends $x$ to
$b$ with $b_i=x_i+t\,u_i$, where $t=(r-\tot(x))/6$; moreover $\tot(b)=r-t$.
\item[\textup{(ii)}] $X_1(\Delta)=P$, and $b_1,\dots,b_4,1-\tot(b)$ are the
barycentric coordinates of $X_1(b)$ with respect to $e_1,\dots,e_4,-u$. Hence
$y\in P$ if and only if $t\ge0$ and $b_i\ge0$ for all $i$, where $b,t$ are
computed from $y$ as in~\textup{(i)} with $r=1$.
\item[\textup{(iii)}] For $r>0$, $X_r(b)=r\,X_1(b/r)$.
\end{enumerate}
\end{lemma}

\begin{proof}
(i) If $x=Mb-ru$ then $\tot(x)=6\tot(b)-5r$, so $t=r-\tot(b)$ and
$x_i+tu_i=b_i+\tot(b)u_i-ru_i+(r-\tot(b))u_i=b_i$. (ii)
$X_1(b)=\sum_ib_ie_i+(1-\tot(b))(-u)$, a convex combination of the vertices
exactly when $b\in\Delta$. (iii) $rX_1(b/r)=Mb-ru$.
\end{proof}

\begin{lemma}\label{lem:points}
The origin is an interior point of $P$, and
$P\cap\Z^4=\{0,-u,e_1,e_2,e_3,e_4\}$. No subset of $P\cap\Z^4$ with at least
two elements has a lattice point as its centroid.
\end{lemma}

\begin{proof}
For $y=0$ and $r=1$ the coordinates of Lemma~\ref{lem:coord} are
$(b,t)=\frac16(2,1,1,1;1)$, all positive. They depend continuously on $y$, so
they stay positive on a neighbourhood of $0$, which lies in $P$.

Let $z\in P\cap\Z^4$ and $s=\tot(z)\in\Z$. By Lemma~\ref{lem:coord}(ii), $s\le1$,
$6z_1+2(1-s)\ge0$ and $6z_i+(1-s)\ge0$ for $i=2,3,4$. Adding the four
inequalities gives $6s+5(1-s)\ge0$, so $s\ge-5$; checking the seven values
$s=-5,\dots,1$ leaves exactly the six points listed, and all of them lie in
$P$.

Let $E\subseteq P\cap\Z^4$ with $|E|\ge2$ have an integral centroid. The
second, third and fourth coordinate sums of $E$ lie in $\{-1,0,1\}$ and are
divisible by $|E|$, so they vanish; hence each of $e_2,e_3,e_4$ is in $E$ if
and only if $-u$ is. If $-u\notin E$, then $E=\{0,e_1\}$, with first
coordinate sum $1$. If $-u\in E$, then $|E|\in\{4,5,6\}$ and the first
coordinate sum is $-2$ or $-1$. In both cases the first coordinate sum is not
divisible by $|E|$.
\end{proof}

\noindent Lean proves $P\cap\Z^4\subseteq\{0,-u,e_1,\dots,e_4\}$ with the same
case analysis, and checks the $64$ subsets of this six-point set by kernel
computation (Section~\ref{sec:lean}). The five vertices lie on the boundary,
so the origin is the only interior lattice point of $P$; this is used only in
Section~\ref{sec:remarks} and is not stated in Lean.

\section{The covering radius}\label{sec:radius}

\begin{lemma}\label{lem:residue}
For $f\in\Z^4$ and $k\in\Z$ let
$\sigma(f,k)=\sum_{i=1}^4\bigl((f_i+ku_i)\bmod 6\bigr)$, with residues in
$\{0,\dots,5\}$.
\begin{enumerate}
\item[\textup{(i)}] For every $f\in\Z^4$ there is $k\in\{0,\dots,5\}$ with
$\sigma(f,k)\le8$.
\item[\textup{(ii)}] For $p=(1,1,2,4)$ and every $k\in\Z$, $\sigma(p,k)\ge8$;
for $k\equiv0,1,2,3,4,5\pmod 6$ the values are $8,13,12,11,10,9$.
\end{enumerate}
\end{lemma}

\begin{proof}
$\sigma(f,k)$ depends only on $f\bmod 6$ and $k\bmod6$, so (i) is a finite
statement about $6^4=1296$ residue vectors and six values of~$k$. It was
checked exhaustively; in Lean by kernel evaluation (the private theorem \texttt{box} in
\texttt{Residues.lean}),
and independently by the following Python lines, which print~$8$:
{\footnotesize
\begin{verbatim}
from itertools import product
u = (2, 1, 1, 1)
score = lambda f, k: sum((a + k*w) % 6 for a, w in zip(f, u))
print(max(min(score(f, k) for k in range(6)) for f in product(range(6), repeat=4)))
\end{verbatim}}
\noindent (ii) is a direct computation.
\end{proof}

\begin{proposition}\label{prop:upper}
Every $x\in\R^4$ can be written as $x=z+2y$ with $z\in\Z^4$ and $y\in P$.
\end{proposition}

\begin{proof}
By Lemma~\ref{lem:coord}(i), $x=X_2(a)$ for some $a\in\R^4$. Let
$f_i=\lfloor6a_i\rfloor$, choose $k$ by Lemma~\ref{lem:residue}(i), and put
$t_i=\lfloor(f_i+ku_i)/6\rfloor$ and $b_i=a_i+ku_i/6-t_i$. Then
\[
6b_i=(6a_i-f_i)+\bigl((f_i+ku_i)\bmod6\bigr),
\]
where the first term lies in $[0,1)$. Hence $b_i\ge0$ and
$6\tot(b)<4+\sigma(f,k)\le12$, so $b/2\in\Delta$. Let
$z=t+(\tot(t)-k)u\in\Z^4$. Since $Mt=t+\tot(t)u$ and $Mu=6u$,
\[
X_2(a)-X_2(b)=M(a-b)=M\bigl(t-\tfrac k6u\bigr)=t+\tot(t)u-ku=z .
\]
By Lemma~\ref{lem:coord}(iii) and~(ii), $X_2(b)=2X_1(b/2)=2y$ with $y\in P$.
\end{proof}

\begin{proposition}\label{prop:lower}
Let $0\le r<2$, $p=(1,1,2,4)$ and $a_i=(p_i+r/2)/6$. Then $X_r(a)$ is not of
the form $z+ry$ with $z\in\Z^4$ and $y\in P$.
\end{proposition}

\begin{proof}
Suppose $X_r(a)=z+ry$. Write $y=X_1(c)$ with $c\in\Delta$
(Lemma~\ref{lem:coord}(ii)). Then $Ma-ru=z+rMc-ru$, so $M(a-rc)=z$. Since
$M^{-1}=I-\frac16u\one^{\mathsf T}$, this gives
$a_i-rc_i=z_i-\tot(z)u_i/6$. With $k=\tot(z)$ and $n_i=p_i+ku_i-6z_i\in\Z$ we get
\[
6rc_i=n_i+r/2 .
\]
The left side is nonnegative and $r/2<1$, so $n_i\ge0$; as
$n_i\equiv p_i+ku_i\pmod6$, also $n_i\ge(p_i+ku_i)\bmod6$. By
Lemma~\ref{lem:residue}(ii), $\tot(n)\ge8$, and therefore
\[
6r\ \ge\ 6r\,\tot(c)\ =\ \tot(n)+2r\ \ge\ 8+2r,
\]
that is, $r\ge2$, a contradiction.
\end{proof}

\noindent Propositions~\ref{prop:upper} and~\ref{prop:lower} give
Theorem~\ref{thm:main}(b): $2$ is the least $r\ge0$ such that
$rP+\Z^4=\R^4$.

\begin{remark}\label{rem:cayley}
By Lemmas~B.1 and~B.2 of~\cite{css}, taking $v_0=-u$ and the volume vector
$(1;2,1,1,1)$ of the origin (six times its barycentric coordinates),
$\mu(P)=(\delta+4)/6$, where $\delta$ is the directed diameter of the Cayley
digraph of $(\Z/6)^4/\langle(2,1,1,1)\rangle$ with respect to the standard
generators. A breadth-first search over its
$216$ vertices gives $\delta=8$, in agreement with $\mu(P)=2$; this computation
is not part of the Lean proof. The search for problems also produced a
different written proof of $\mu(P)\le2$, which uses a sawtooth function and
the Cauchy--Schwarz inequality; it is not used here.
\end{remark}

\section{Volume and the equality list}\label{sec:list}

\paragraph{Volume.} The affine map $b\mapsto Mb-u$ sends $0$ to $-u$ and $e_i$
to $e_i$, hence $\conv\{0,e_1,\dots,e_4\}$ onto $P$. Its linear part has
determinant~$6$, so $\vol(P)=6\vol(\Delta)$. $\Delta$ is compact, so $\vol(\Delta)<\infty$,
and $\vol(\Delta)=\vol(P)/6>0$ because $P$ has interior points; hence $\Vol(P)=6$. This is
Theorem~\ref{thm:main}(c).

\begin{definition}\label{def:G}
Let $\G$ be the smallest family of subsets of $\R^4$ such that
\begin{enumerate}
\item[(L)] $\conv(s)\in\G$ for every finite set $s\subseteq\Z^4$ with
$|s|\ge2$ and $\sum_{z\in s}z=0$;
\item[(T)] if $K\in\G$ and $z\in\Z^4$, then $z+K\in\G$;
\item[(J)] if $K,L\in\G$, then $\conv(K\cup L)\in\G$.
\end{enumerate}
\end{definition}

\noindent Rule~(J) imposes no condition on subspaces, lattices or the
position of the origin, so $\G$ is designed to be larger than the list of
Conjecture~A (Lemma~\ref{lem:G}).

\begin{lemma}[not in Lean]\label{lem:G}
Every polytope in $\R^4$ that is obtained from simplices $S(\one_l)$ by direct
sums and/or lattice translations, as in~\cite{css}, belongs to $\G$, and so
does its image under any unimodular transformation.
\end{lemma}

\begin{proof}
In the construction of~\cite{css} each intermediate polytope $Q$ lies in a
real vector space $V_Q$ with a lattice $\Lambda_Q$: a simplex $S(\one_l)$
lies in $\R^{l-1}$ with $\Z^{l-1}$; a lattice translate $\lambda+K$
($\lambda\in\Lambda_K$) lies in $(V_K,\Lambda_K)$; and $K\oplus L$ lies in
$V_K\oplus V_L$ with $\Lambda_K\oplus\Lambda_L$. The final polytope becomes a
polytope in $\R^4$ when $(V_Q,\Lambda_Q)$ is identified with $(\R^4,\Z^4)$ by
a linear isomorphism. We show by induction on the construction: \emph{for
every injective linear map $\phi\colon V_Q\to\R^4$ with
$\phi(\Lambda_Q)\subseteq\Z^4$, we have $\phi(Q)\in\G$.} If $Q=S(\one_l)$
with $l\ge2$, put $b_j=\phi(e_j)$. Then
$\phi(Q)=\conv\{-\sum_jb_j,b_1,\dots,b_{l-1}\}$, and these $l$ points are
distinct lattice points (the $b_j$ are linearly independent) with sum zero,
so $\phi(Q)\in\G$ by~(L). The point $S(\one_1)=\{0\}$ satisfies
$K\oplus\{0\}=K$ and can be omitted. If $Q=\lambda+K$ with
$\lambda\in\Lambda_K$, then $\phi(Q)=\phi(\lambda)+\phi(K)$ with
$\phi(\lambda)\in\Z^4$, which is~(T). If $Q=K\oplus L$, then
$Q=\conv(K\cup L)$: every point $\lambda x+(1-\lambda)y$ lies in the hull,
and conversely a convex combination of points of $K$ and of $L$ can be
grouped as $\lambda x+(1-\lambda)y$ with $x\in K$, $y\in L$, by convexity of
$K$ and $L$. This uses neither the origin nor the lattices, so it also holds
when the origin lies on the boundary of a factor. Hence
$\phi(Q)=\conv(\phi(K)\cup\phi(L))$. The restrictions of $\phi$ to $V_K$ and
$V_L$ are injective and map $\Lambda_K$ and $\Lambda_L$ into $\Z^4$, so the
induction hypothesis and~(J) apply. Taking for $\phi$ the identifying
isomorphism gives the first claim. For the second, a matrix $U\in GL_4(\Z)$ maps sets as in~(L) to
sets as in~(L), and commutes with convex hulls and maps $z+K$ to $Uz+UK$; so
$U\G\subseteq\G$ by induction, and a final translation by a lattice vector is
rule~(T).
\end{proof}

\begin{proposition}\label{prop:outside}
$P\notin\G$.
\end{proposition}

\begin{proof}
By induction on the generation of $K\in\G$, there are a finite set
$s\subseteq\Z^4$ with $|s|\ge2$ and $\sum_{z\in s}z=0$, and a vector
$t\in\Z^4$, such that $t+s\subseteq K$: for~(L) take $t=0$; for~(T) replace
$t$ by $z+t$; for~(J) keep the set found for $K$. If $P\in\G$, then $t+s$ is a
set of $|s|\ge2$ distinct lattice points of $P$ whose centroid is the lattice
point $t$. This contradicts Lemma~\ref{lem:points}.
\end{proof}

\noindent This proves Theorem~\ref{thm:main}(d), and with
Lemma~\ref{lem:G}, Corollary~\ref{cor:A}. The exclusion does not use the
volume.

\section{Further remarks (not in Lean)}\label{sec:remarks}

The statements in this section are written arguments. None of them is part
of the Lean statements.

\begin{remark}[arbitrary translations]\label{rem:real}
If ``translations'' in Conjecture~A are read as arbitrary real translations,
while direct sums keep the conditions of~\cite[Section~2.1]{css} (both factors
contain the origin, and the lattice is the direct sum of the factor lattices),
then $P$ is still not in the list, by the following volume argument, which is
not formalized. Translations preserve volume. For a direct sum, triangulate
each factor into simplices with apex at the origin over a triangulation of
the facets of the factor that do not contain the origin (only full-dimensional pieces are kept, so this also
works when the origin is on the boundary). The direct sum of two such pieces
is a simplex whose edge matrix is block diagonal in a basis of the direct sum
lattice, so normalized volumes multiply, and these simplices tile $K\oplus L$
with disjoint interiors. Hence normalized volume is multiplicative on direct
sums. Since $\Vol(S(\one_{j+1}))=j+1$, the polytopes of the list in dimension
four have normalized volume $5,8,9,12$ or $16$, for the dimension partitions
$4$, $3+1$, $2+2$, $2+1+1$, $1+1+1+1$. None of these is $6=\Vol(P)$.
\end{remark}

\begin{remark}[Conjecture C and Corollary~4.3 of~\cite{css}]\label{rem:C}
For a $d$-simplex $S$ with the origin in its interior and rational vertex
directions, Conjecture~C of~\cite{css} states $\mu(S)\le\frac12\,\mathrm{Surf}(S)/\Vol(S)$,
which by~\cite[Lemma~4.1]{css} is equivalent to
$\mu(S)\le\frac12\sum_i1/\ell_i$, where $\ell_i$ is the lattice length of the
segment from the vertex $v_i$ through the origin to the opposite facet.
Corollary~4.3 of~\cite{css} is stated as ``Conjecture~C $\Rightarrow$
Conjecture~A''; its proof derives the inequality $\mu(S)\le d/2$ for lattice
simplices and does not address the equality clause. For $P$ every vertex is
primitive, so $\ell_i=1/(1-\lambda_i)$, where $\lambda=\frac16(2,1,1,1,1)$ are
the barycentric coordinates of the origin, and
$\frac12\sum_i(1-\lambda_i)=\frac12(5-1)=2=\mu(P)$. So $P$ satisfies the
inequality of Conjecture~C with equality and does not refute it. So Conjecture~C,
if true, cannot imply the equality clause of Conjecture~A, which fails for
$P$; Corollary~4.3 holds for the inequality, which is what its proof
shows.
\end{remark}

\begin{remark}[Conjecture D of~\cite{css}]\label{rem:D}
Conjecture~D of~\cite{css} states that a lattice $d$-polytope with $k$
interior lattice points has $\mu(P)\le\frac{d-1}2+\frac1{k+1}$, and that for
$k=1$ equality holds if and only if $P$ is obtained by direct sums and/or
translations of simplices $S(\one_l)$. Our $P$ has exactly one interior
lattice point (Section~\ref{sec:coord}) and $\mu(P)=2=\frac32+\frac12$, so it
also contradicts the equality clause of Conjecture~D for $k=1$, with the same
reading of ``translations''. The Lean lemma \texttt{Exclusion.lattice\_point}
determines the lattice points of $P$, but Conjecture~D is not a Lean target.
\end{remark}

\begin{remark}[the simplices $S(\omega)$]\label{rem:somega}
Theorem~1.4 of~\cite{css} gives $\mu(S(\omega))$ in closed form. Suppose a
unimodular transformation $\phi$ mapped some $S(\omega)$ onto $P$. The origin is
an interior lattice point of $S(\omega)$, so $\phi(0)$ is an interior lattice
point of $P$, hence $\phi(0)=0$ and $\phi$ is linear. The vertices
$\omega_ie_i$ and $-\omega_0\one_d$ are mapped to vertices of $P$, which are
lattice points; so they are lattice points, and the $\omega_i$ are positive
integers. Since the vertices of $P$ are primitive, so are their preimages, and
all $\omega_i=1$. But $S(\one_5)$ has normalized volume $5\ne6$. So $P$ is not
of the form $S(\omega)$ up to unimodular equivalence, and Theorem~1.4
of~\cite{css} does not give $\mu(P)$.
\end{remark}

\begin{remark}[other written material]
The search for problems in which this example was found also produced written
arguments for related statements, among them variants of the example in every
dimension $d\ge4$ and an example concerning the equality clause of
Conjecture~D for $k\ge2$. These are written arguments, have not been verified
in Lean or reviewed for this paper, and are not claims of this paper.
\end{remark}

{\sloppy
\section{Formal verification}\label{sec:lean}

Theorem~\ref{thm:main} is formally verified in Lean~4 (version 4.34.1) with
Mathlib~\cite{mathlib} (commit \texttt{d13f23b7}). Lemma~\ref{lem:G},
Corollary~\ref{cor:A} beyond Theorem~\ref{thm:main}, and
Section~\ref{sec:remarks} are not in Lean. The $9$ files are in
version~\repoversion{} of the repository~\cite{repo}, directory
\texttt{lean/Research/TaskR/CSS/}.

\paragraph{The statements.}
The statement file \texttt{Challenge.lean} contains only definitions and
statements, no proofs. In ASCII transliteration (the docstrings and the
\texttt{import}, \texttt{namespace} and \texttt{end} lines omitted; the file
itself uses Unicode, e.g.\ $\mathbb R$, $\to$, $\forall$, $\exists$, $\in$,
$\le$, $\wedge$, $\neg$, $\cup$, $\sum$):
{\footnotesize
\begin{verbatim}
open Set MeasureTheory
open scoped BigOperators ENNReal
abbrev Space := Fin 4 -> Real
abbrev Lattice := Fin 4 -> Int
def embed (z : Lattice) : Space := fun i => (z i : Real)
def weight : Lattice := ![2, 1, 1, 1]
def e (i : Fin 4) : Lattice := Pi.single i 1
def vertices : Set Space := insert (embed (-weight)) (range fun i => embed (e i))
noncomputable def P : Set Space := convexHull Real vertices
noncomputable def unitSimplex : Set Space :=
  convexHull Real (insert 0 (range fun i => embed (e i)))
noncomputable def normalizedVolume (K : Set Space) : ENNReal := volume K / volume unitSimplex
def Covers (K : Set Space) (r : Real) : Prop :=
  forall x : Space, exists z : Lattice, exists y in K, x = embed z + r * y
inductive Generated : Set Space -> Prop
  | leaf (s : Finset Lattice) (hcard : 2 <= s.card) (hsum : sum z in s, z = 0) :
      Generated (convexHull Real (embed '' (s : Set Lattice)))
  | translate {K : Set Space} (h : Generated K) (z : Lattice) :
      Generated ((fun x => embed z + x) '' K)
  | join {K L : Set Space} (hK : Generated K) (hL : Generated L) :
      Generated (convexHull Real (K union L))
def NonHollow : Prop := (0 : Space) in interior P
def Radius : Prop := IsLeast {r : Real | 0 <= r /\ Covers P r} 2
def Volume : Prop := normalizedVolume P = 6
def OutsideList : Prop := not Generated P
\end{verbatim}}
\noindent Here \texttt{ENNReal} is the type of extended nonnegative reals,
$[0,\infty]$, written in the file with the notation
$\mathbb R{\ge}0\infty$, which needs the scope
\texttt{ENNReal} opened in the second line; \texttt{r * y} stands for the
scalar multiple $r\bullet y$; \texttt{embed '' s} is the image of $s$ under \texttt{embed}
(the coercion arrow before \texttt{s} is omitted above); \texttt{volume} is
Lebesgue measure on $\R^4$; and \texttt{IsLeast A 2} says that $2\in A$ and
$2\le r$ for all $r\in A$. \texttt{P} is the convex hull of $-u$ and
$e_1,\dots,e_4$. \texttt{NonHollow} is Theorem~\ref{thm:main}(a); it implies
that $P$ is $4$-dimensional and non-hollow. \texttt{Covers P r} says that
$rP+\Z^4=\R^4$, so \texttt{Radius} says that the minimum in the definition of
$\mu(P)$ exists and equals~$2$, which is Theorem~\ref{thm:main}(b).
\texttt{normalizedVolume} divides by the volume of the unimodular simplex
$\conv\{0,e_1,\dots,e_4\}$; \texttt{Volume} is Theorem~\ref{thm:main}(c).
\texttt{Generated} is the family $\G$ of Definition~\ref{def:G}, and
\texttt{OutsideList} is Theorem~\ref{thm:main}(d).

The final reviewer (see the last section) checked the following points. (i)~The
covering radius is the minimum of~\cite{css}, not a finite proxy, and the
reviewer restated \texttt{Radius} with its own vertex list and an explicit
$\exists z\,\exists y$ and derived it from the checked theorem. (ii)~The
quotient in \texttt{normalizedVolume} could be degenerate only if the volume of
the unit simplex were $0$ or $\infty$; the proof excludes both. (iii)~The
family \texttt{Generated} is not empty (it contains the segment
$[-e_1,e_1]$), so \texttt{OutsideList} is not vacuous; it contains the list
of Conjecture~A with lattice translations (Lemma~\ref{lem:G}), so
\texttt{OutsideList} is at least as strong as ``$P$ is not in that list''.
(iv)~\texttt{Generated} uses lattice translations only, matching p.~2
of~\cite{css}; real translations are covered only by
Remark~\ref{rem:real}.

\paragraph{Review and freezing.}
The statement went through two versions. Version~1 of \texttt{Challenge.lean}
(SHA-256 \texttt{7fff04f3\ldots}) was checked against~\cite{css}, together
with the reading of the source, by a Codex agent that had written neither the
statement nor any proof and had not taken part in the earlier proofs of this
example; it approved the file on 20 November 2026 at 19:31:21~UTC. The file
had not been compiled before this approval. Its compilation (log
19:31:49--19:32:19~UTC) failed with a parse error, because the notation for
extended nonnegative reals needs the \texttt{ENNReal} scope. Version~2
(SHA-256 \texttt{1dbf0ce8\ldots}) differs from version~1 only in line~6,
where \texttt{open scoped BigOperators} became \texttt{open scoped
BigOperators ENNReal}; all definitions and the four target lines are
byte-identical. Version~2 compiled without diagnostics (19:32:27--19:32:45~UTC).
A second Codex agent, which had not written the statement or the proofs for
this problem (it was writing the Lean proof of another problem in the same
task), compared the two versions and approved version~2 at 19:33:19~UTC, and
version~2 was frozen by recording its SHA-256 hash at 19:33:34~UTC. Between the
approval of version~1 and the freeze, two proof files were created (file
creation times): \texttt{Geometry.lean} at 19:31:49~UTC, which was revised
after the freeze (last change 19:34:14~UTC), and \texttt{Residues.lean} at
19:32:14~UTC, which was not changed afterwards and is the checked file.
\texttt{Residues.lean} uses none of the definitions of \texttt{Challenge.lean}. According to the development logs, no
proof file was compiled before 19:33:41~UTC, after the freeze. The statement file was not changed after
the freeze. The final reviewer diffed the two versions and judged the
difference immaterial: opening the scope adds only notation, and the
mathematical content that the proofs target was approved before any proof
file existed. A process rule was recorded for later tasks: a statement must
compile without diagnostics before it is approved for freezing.

\paragraph{The proofs.}
\texttt{Proof.lean} defines \texttt{Proof.NonHollow} as
\texttt{Geometry.nonHollow}, \texttt{Proof.Radius} as \texttt{radius}
(file \texttt{Radius.lean}, namespace \texttt{TaskR.CSS.Proof}) and
\texttt{Proof.Volume} as \texttt{VolumeCalc.volume\_\allowbreak
from\_\allowbreak nonHollow} applied to \texttt{Proof.NonHollow};
\texttt{Proof.OutsideList} is proved in \texttt{Exclusion.lean}.
Table~\ref{tab:lean} lists the Lean counterparts of the steps of the written
proof (names relative to the namespace \texttt{TaskR.CSS}; ``private''
marks private declarations).

\begin{table}[ht]
\centering\footnotesize
\renewcommand{\arraystretch}{1.05}
\begin{tabular}{@{}>{\raggedright\arraybackslash}p{0.42\textwidth}>{\raggedright\arraybackslash}p{0.54\textwidth}@{}}
\hline
Step of the written proof & Lean (file) \\
\hline
$\tot(x)$, $t$, $b$ and $X_r$ of Lemma~\ref{lem:coord}(i) &
\texttt{Geometry.total}, \texttt{slack}, \texttt{coeff}, \texttt{mix};
\texttt{coeff\_mix}, \texttt{mix\_coeff}, \texttt{total\_coeff},
\texttt{total\_mix} (\texttt{Geometry}) \\
Lemma~\ref{lem:coord}(ii): $X_1(\Delta)\subseteq P$ and the converse &
\texttt{Geometry.mem\_mix}, \texttt{Geometry.constraints},
\texttt{Geometry.mem\_of\_constraints} (\texttt{Geometry}) \\
Lemma~\ref{lem:points}: origin interior &
\texttt{Geometry.nonHollow} (\texttt{Geometry}) \\
Lemma~\ref{lem:points}: lattice points of $P$ &
\texttt{Exclusion.lattice\_point} (\texttt{Exclusion}), from
\texttt{Exclusion.finite\_classification} and the private
\texttt{coord\_cases} (\texttt{ExclusionFinite}) \\
Lemma~\ref{lem:points}: no integral centroid ($64$ subsets) &
private \texttt{nonintegral\_subsets}, \texttt{Exclusion.points}
(\texttt{ExclusionFinite}) \\
Lemma~\ref{lem:residue} &
\texttt{Residues.score}, private \texttt{box}, \texttt{Residues.upper},
\texttt{Residues.lower} (\texttt{Residues}) \\
Proposition~\ref{prop:upper} &
private \texttt{upper\_coordinates}, \texttt{coversTwo} (\texttt{Radius}) \\
Proposition~\ref{prop:lower} &
private \texttt{base}, \texttt{lower\_coordinates} (\texttt{Radius}) \\
Theorem~\ref{thm:main}(b) & \texttt{Proof.radius} (\texttt{Radius}) \\
Volume: $M$, $\det M=6$, image of $\Delta$ &
private \texttt{edges}, \texttt{det\_edges}, \texttt{placement},
\texttt{image\_body}; \texttt{VolumeCalc.volume\_from\_nonHollow}
(\texttt{Volume}) \\
Proposition~\ref{prop:outside}: induction and centroid &
private \texttt{contains\_leaf} (\texttt{Exclusion}),
\texttt{Exclusion.no\_translated\_leaf} (\texttt{ExclusionFinite}),
\texttt{Proof.OutsideList} (\texttt{Exclusion}) \\
\hline
\end{tabular}
\caption{Lean counterparts of the steps of the proof.}
\label{tab:lean}
\end{table}

\noindent The test file \texttt{CheckCSS.lean} restates the four targets with
the frozen types (\texttt{check\_\allowbreak NonHollow},
\texttt{check\_\allowbreak Radius}, \texttt{check\_\allowbreak Volume},
\texttt{check\_\allowbreak OutsideList}, each of the form
\texttt{theorem check\_X : Challenge.X := Proof.X}). The axioms of these four
theorems and of the four \texttt{Proof.*} theorems, printed during
compilation, are exactly \texttt{propext}, \texttt{Classical.choice} and
\texttt{Quot.sound}. No file uses \texttt{sorry}, \texttt{native\_\allowbreak
decide}, \texttt{bv\_\allowbreak decide} or an \texttt{axiom} declaration.
The two finite checks use \texttt{decide +kernel} and are checked by the
kernel: \texttt{box} (\texttt{Residues.lean}; $6^4=1296$ residue vectors, six values of
$k$; the only place with raised \texttt{maxRecDepth} and
\texttt{maxHeartbeats}) and \texttt{nonintegral\_\allowbreak subsets}
($64$ subsets). The other uses of \texttt{decide} prove facts such as
$(6:\Z)\ne0$.

\paragraph{Checks.}
The $9$ files ($6$ proof modules, \texttt{Challenge.lean},
\texttt{Proof.lean} and \texttt{CheckCSS.lean}) were compiled module by
module in two new empty directories on the author's Mac, with the pinned
toolchain and with warnings treated as errors, on 20 November 2026
(19:35:47--19:40:41~UTC), and every module was replayed with
\texttt{leanchecker}; each directory also ran one build of the whole package.
All $38$ commands ($19$ per directory) ended with exit code~$0$, the
\texttt{leanchecker} runs printed nothing, and the two directories produced
identical \texttt{.olean} files and axiom lists. The final reviewer (see the
next section) compiled the $9$ files again, together with a short audit file
of its own, in a third new directory on the same Mac on 20 November 2026
(between 19:48 and 19:52~UTC, the times of its lock on the machine), and
replayed every module with \texttt{leanchecker}; all runs ended with exit
code~$0$, with no warnings and no \texttt{leanchecker} output, and the axiom
lists were the same. Mathlib and its dependencies were taken from a shared
prebuilt cache at the pinned commits and were not rechecked; the
\texttt{leanchecker} replays cover the $9$ project modules. In the repository
the files have module names under \texttt{Research.TaskR.CSS} and differ from
the checked files only in their \texttt{import} lines
(\texttt{Challenge.lean} is unchanged). The repository versions were made by a script that changed only the local \texttt{import} lines, and on 20 November 2026 (19:55--20:07~UTC) they were compiled and replayed with \texttt{leanchecker} again, module by module, in two new empty directories on the author's Mac, one after the other, together with the files of three other papers of version~\repoversion{} ($39$ modules in all); there were no warnings, and the two directories produced identical \texttt{.olean} files and the same axiom lists. In the whole-project rebuild of version~\repoversion{} of the repository (on the author's Windows computer, in two passes on 20 and 21 November 2026), these files were compiled and replayed with \texttt{leanchecker} again in pass~2 (21 November 2026), which built the files not included in pass~1 on top of its results; no module needed a second attempt (details in the repository). The
formalization has not been checked with comparator or with an independent
kernel.

\par}

\section*{How this paper was produced}
This work was carried out from 20 to 22 November 2026 (UTC) in a workflow directed by
the author, using OpenAI Codex with the model GPT-6.2-Astra (reasoning level
ultra) and Anthropic Claude Code with the model Claude Opus~5.5. The author
asked for a search for hard and substantial open problems, and Claude Code
passed this request to Codex. Between 18:36 and 18:54~UTC, in a separate Codex
session, Codex screened $78$ candidate problems and prepared detailed notes on
eight of them. The example of this paper was found and written up in this
search, checked by exact computation, and read by another Codex agent (same
model) working in a parallel lane of the search, which gave an independent
written proof of the upper bound. After reading the report, the author chose
this problem and one other. From 19:29~UTC (time of the task description),
Codex agents (same model) checked the reading of the source, searched again for
prior work, wrote the statement file, the Lean proof and the written proof,
and ran the checks; the review and freezing of the statement, with its two
versions, are described in Section~\ref{sec:lean}. According to the file
timestamps and the development logs, two proof files were created between the
approval of version~1 and the freeze of version~2 (one of them in its final
form), all other proof files were
created after the freeze, the six proof modules were compiled during
development between 19:33:41 and 19:35:33~UTC, and \texttt{Proof.lean} and
\texttt{CheckCSS.lean} were first compiled in the acceptance run; they follow the written proof prepared in the search for
problems. After the checks, an independent Claude Code agent (model Claude
Opus~5.5) that had written none of the files and had taken no part in the
search or the formalization gave the final review. It read the source itself
(arXiv versions~1 and~2 and the journal text), compared the statement file
with it, read the written proof, the proof files, the freeze records and the
raw logs of the checks, recomputed the lattice points, facet normals
and determinant in Python and checked the covering radius numerically on
fine grids, and rebuilt the formalization as
described above. Its decision was APPROVE\_FORMAL, with four minor remarks on
the written account, all taken into account here: the reading of
Corollary~4.3 and Conjecture~D (Remarks~\ref{rem:C} and~\ref{rem:D}), the
lattice translations (Remark~\ref{rem:real}), the earlier appearances of the
simplex and the simplices $S(\omega)$ (Prior work and
Remark~\ref{rem:somega}), and the freezing procedure
(Section~\ref{sec:lean}). Claude Code (Claude Opus~5.5) checked the sources,
repeated the finite computations of Lemmas~\ref{lem:points}
and~\ref{lem:residue} and Remark~\ref{rem:cayley} in Python, and wrote this
text from Codex's manuscript and the Lean files.
An independent Claude Code agent (model Claude Opus~5.5) that had written none of the files reviewed this paper, and another reviewed the release; their corrections were applied. Throughout, ``we'' refers to this workflow; the reading and
searching described in the paper were done by the agents. Codex agents
searched for prior work on 20 November 2026, in the search for problems and
again before the statement was frozen, and the final reviewer searched again;
they found no earlier statement of the covering radius of $P$ and no
counterexample to the equality clause of Conjecture~A of~\cite{css} (details
in the repository). Claude Code searched for works citing~\cite{css} on 20 November 2026 (23:17--23:43~UTC) and again, together with the arXiv, Semantic Scholar, Crossref, DataCite, Zenodo, MathDB, GitHub and repositories of AI-assisted mathematics, on 21 November 2026 (17:54--18:21~UTC), and found no new related work; on 21 November the citation index OpenAlex refused to list the works citing~\cite{css} (rate limit), and Semantic Scholar listed the same ten citing works as before. Claude Code searched again on 22 November 2026 (05:39--05:51~UTC), when the OpenAlex lists of citing works refused on 21 November were retrieved, and found nothing new. This reports the coverage of the
searches, not a guarantee that they were complete. The author is not a
professional mathematician. The author has read this paper and, using a side-by-side table, compared its Lean statements with the main theorem, Theorem~\ref{thm:main}. No human mathematician has reviewed
this paper. The correctness of the results rests on the Lean proofs and on the
agreement of the Lean statements with Conjecture~A of~\cite{css}, together with
the argument of Lemma~\ref{lem:G}; the written proofs are given for the
reader. AI systems are not authors of this paper; the author takes full
responsibility for its content.

\medskip
\noindent Thanks are due to Giulia Codenotti, Francisco Santos and Matthias
Schymura for the conjecture.

{\small
\begin{thebibliography}{9}
\setlength{\itemsep}{1pt}
\bibitem{css} G.~Codenotti, F.~Santos and M.~Schymura, The covering radius and
a discrete surface area for non-hollow simplices, \emph{Discrete Comput.
Geom.} 67 (2022), no.~1, 65--111,
\url{https://doi.org/10.1007/s00454-021-00330-3}; arXiv:1903.02866v2 (2021).
\bibitem{cfk} G.~Codenotti, A.~Freyer and K.~Krivokuća, Minimal covering
bodies and Brunn--Minkowski type inequalities for the covering radius,
arXiv:2606.14442, 2026.
\bibitem{repo} H.~Dong, \emph{ai4math-results}, repository,
\url{https://github.com/Horace-Maxwell/ai4math-results}, version~\repoversion{}
(2026); archived at Zenodo under the concept DOI
\url{https://doi.org/10.5281/zenodo.22970254}.
\bibitem{kas} A.~M. Kasprzyk, \emph{Toric Fano varieties and convex
polytopes}, PhD thesis, University of Bath, 2006,
\url{https://purehost.bath.ac.uk/ws/portalfiles/portal/187944797/kasprzyk_20thesis_202006.pdf}.
\bibitem{kr} K.~Krivokuća, Upper bounds on covering minima of convex bodies,
arXiv:2601.15173v1, 2026.
\bibitem{lee} T.-J. Lee, On the hyperplane conjecture, arXiv:2610.26186v2,
2026 (Appendix~C; the same appendix is in v1).
\bibitem{mathlib}
The mathlib Community, The Lean mathematical library, in \emph{Proc.\ 9th ACM
SIGPLAN Int.\ Conf.\ Certified Programs and Proofs} (CPP '20), ACM, 2020,
\url{https://doi.org/10.1145/3372885.3373824}.
\end{thebibliography}}
\end{document}
